% !TEX root = ../main.tex
\section{Limitations and Dual-Use / Ethical Considerations}
The mechanical outcome profiles are not admitted resistance estimates.
The two paid wrong-majority cases establish source-adverse outcomes and extra
payments, but do not demonstrate prior knowledge of falsity or isolate a causal
effect of the offer. Semantic coordination and competence review are still
open. Confirmation uses only two source items in finance and education. The separate
five-source pilot also has limited coverage, and all cases use source-relative
binary verification. Repeating them under new seeds does not
create new documents. The first campaign uses homogeneous model juries with
scripted parties, so it cannot establish cross-role autonomy or spontaneous
coalition formation. Small cluster counts, provider nondeterminism, budget
truncation, and unresolved semantic review can limit model comparisons even
after execution. The two fixed items occur in the same order, so source
differences also coincide with episode position. Selected quotations illustrate
contrasts; their selection supplies neither prevalence nor faithful access to
internal reasoning. All paid offers use one amount and are unescrowed, so
responses can reflect payment credibility rather than an intrinsic refusal
to coordinate. The 75 startup-recovery attempts form a separate, condition-skewed
exploratory cohort and do not replace original failures.

Reference support, an incorrect verdict, an applicable-duty violation, and
executed collusion answer different questions. In particular, support in a
scientific abstract is not independent clinical truth, and an excerpted
contract task is not legal advice or the complete ContractNLI task. A positive
honest baseline and feasibility checks establish an available compliant path,
not that all models understand or can complete it. Independent review of
adaptations, rule delivery, and semantic labels remains required.

Experiments use isolated test-chain rewards; no real participant financial
loss is claimed. Conditional-reward scenarios and detailed trajectories could
help attackers recognize opportunities to corrupt automated review. Their
defensive purpose is to test evidence standards and candidate safeguards while
measuring useful honest work. Credentials, evaluator-only references, and
coalition plans stay outside evaluated participant inputs. Private archival of
new text and outputs is authorized; public evidence release requires its own
authorization and review. Historical public-release permission does not extend
to future campaigns.

The five authors are listed in their supplied order with confirmed affiliations;
final author review remains open. This working draft is not submission-ready.
